\section{Discussion}
\vspace{-0.3cm}
In this work, we introduced a new method for achieving robust group-invariance in group-equivariant convolutional neural networks. Our approach, the \textit{$G$-TC Layer}, is built on the \textit{triple correlation} on groups, the lowest-degree polynomial that is a complete group-invariant map \cite{kakarala1992triple, sturmfels2008algorithms}. Our method inherits its completeness, which provides measurable gains in robustness and classification performance as compared to the ubiquitous Max $G$-Pooling.

This improved robustness comes at a cost: the $G$-TC Layer increases the dimension of the output of a $G$-Convolutional layer from $G$ to $\frac{|G|(|G|+1)}{2}$. While the dimension of the discretized groups used in $G$-CNNs is typically small, this increase in computational cost may nonetheless deter practitioners from its use. However, there is a path to further reduction in computational complexity provided that we consider its spectral dual: the bispectrum. In \cite{kondor2008thesis}, an algorithm is provided that exploits more subtle symmetries of the bispectrum to demonstrate that only $|G| + 1$ terms are needed to provide a complete signature of signal structure, for the one-dimensional cyclic group. In Appendix \ref{app:reduction}, we extend the computational approach from \cite{kondor2008thesis} to more general groups and provided a path for substantial reduction in the complexity of the $G$-TC Layer, thus expanding its practical utility. Novel mathematical work that grounds our proposed computations in group theory is required to quantify the exact complexity reduction that we provide.


As geometric deep learning is applied to increasingly complex data from the natural sciences \cite{gainza2020deciphering, atz2021geometric, ju2021performance}, we expect robustness to play a critical role in its success. Our work is the first to introduce the general group-invariant triple correlation as a new computational primitive for geometric deep learning. We expect the mathematical foundations and experimental successes that we present here to provide a basis for rethinking the problems of invariance and robustness in deep learning architectures.
